
\section{Introduction}
\label{ch_intro}
Einstein's theory of relativity unifies space and time, yet it retains a structural asymmetry: space is a three-dimensional vector domain with intrinsic rotational freedom, whereas time enters as a one-dimensional direction. Expressing local displacements within their respective subspaces as $\mathbf{x}_{s}=x^{s}\mathbf{e}_{s}$ and $\mathbf{x}_{t}=x^{0}\mathbf{e}_{0}$ makes this asymmetry explicit and motivates the fundamental question: \emph{can time possess a three-dimensional vector structure without sacrificing the four-dimensional results of relativity or introducing the causal pathologies common to multi-time formulations; and, if so, how do the resulting dynamics connect to established physics?}

To address this question, we propose a framework built on the Principle of Relativity and the Principle of Maximal Spatiotemporal Symmetry. The latter posits that the temporal subspace possesses the same three-dimensional vector structure and algebraic status as the spatial subspace by default, unless physical consistency demands otherwise. From these premises, we demonstrate that the resulting framework is consistent and constitutes a single theory unifying special relativity and quantum physics.

Within this construction, temporal-frame compatibility with the relativistic invariant $c$ confines causal temporal displacement to a single stabilized temporal axis, thereby establishing a rigorous arrow of time. This constraint decomposes the six-dimensional spacetime into complementary $4+2$ sectors: a four-dimensional translational spacetime that exactly recovers Einsteinian relativity and an orthogonal two-dimensional temporal plane, termed the Plane-of-Now. The rotational structure of this Plane-of-Now gives rise to phase and spin degrees of freedom, Abelian and non-Abelian gauge structures, as well as mass, Yukawa, and isospin interactions, all without requiring additional independent postulates.

Thus, the Principle of Maximal Spatiotemporal Symmetry enlarges the underlying geometry without disturbing its established relativistic content. The resulting framework provides a common geometric foundation for the principal quantum and field-theoretic structures recovered by the theory while remaining exactly consistent with relativity and four-dimensional translational reality.

\subsection{Related Approaches in the Literature}
Endowing time with additional dimensions is not itself novel \cite{demers1975_vector_time,cole1980_particle_decay,kletetschka2025_three_dimensional_time}. Previous approaches employ temporal-orientation averaging \cite{cole1978_six_dimensional_relativity}, gauge constraints yielding one-time descriptions \cite{bars2000_2t_field_theory,bars_rey2001_sp2r_2t}, independent temporal variables and evolution parameters \cite{dirac1932_relativistic_qm,petrat_tumulka2014_multitime_qft,stueckelberg1941_temps_propre,horwitz_piron1973_relativistic_dynamics}, nonlocal constraints on initial data \cite{craig_weinstein2009_multiple_time}, or symmetry breaking that selects one observable time while leaving other temporal directions hidden \cite{DartoraCabrera2010}, while geometric interpretations of quantum phase have also been developed within four-dimensional spacetime \cite{hestenes2015_spacetime_algebra,hestenes1967_real_spinor_fields}.

The present theory does not use these mechanisms to restrict the temporal geometry. Instead, time remains a three-dimensional on-shell vector subspace with a structure symmetric to that of space. Compatibility with the relativistic invariant follows solely from the Principle of Relativity and the Principle of Maximal Spatiotemporal Symmetry, with equality between the six-dimensional velocity norm and its four-dimensional Lorentzian expression fixing the causal temporal direction. Accordingly, causal translation is confined to one temporal direction while the remaining two remain active as the Plane-of-Now, providing the geometric setting for quantum phenomena without additional causal clocks, compactification, gauge elimination of temporal components, hidden temporal directions, or supplementary geometric or dynamical constraints.

\subsection{Scope and Interpretation}
Before proceeding, two points must be clarified. First, the purpose of this foundational paper is to answer the question posed above and to establish the viability and internal coherence of the framework; it does not aim to exhaust every solution or experimental consequence. Second, this framework is an independent geometric theory, not a reformulation of standard quantum mechanics. No quantum structures or explicit relations to standard quantum mechanics are assumed a priori. The theory should therefore be evaluated strictly within its own postulates and internal consistency.

\subsection*{Notation and Conventions}
Bold calligraphic uppercase symbols denote six-dimensional vectors or full Clifford objects, such as $\boldsymbol{\mathcal{X}}$ and $\boldsymbol{\mathcal{Q}}$, while bold regular uppercase symbols denote four-dimensional vectors or Clifford operators in the translational sector, such as $\mathbf{U}$ and $\mathbf{G}$. Bold lowercase symbols denote objects restricted to a specified subspace, such as $\mathbf{x}_{s}$ and $\mathbf{x}_{t}$. Component notation follows a corresponding distinction: four-dimensional components are written with uppercase symbols and spacetime indices $\mu,\nu\in\{0,1,2,3\}$, as in $X^{\mu}$ and $P_{\mu}$; lowercase scalar components are retained for quantities defined within a specified spatial or temporal subspace, or used as scalar coordinates in the six-dimensional construction, as in $x^{s}$, $x^{t}$, and $x^{\alpha}$. Finally, the notation $O=\rho(\mathbf{O})$ identifies a chosen matrix representation of an abstract operator $\mathbf{O}$.
